% New-round theorem fragment. No numerical all-sign norm bound is asserted.
% Requires amsmath, amssymb, amsthm and proposition/proof environments.
\section{Finite odd multipliers with an identity tail}
Let $P_j^{(d)}$ be the orthogonal projection onto total Gaussian Hermite
degree $j$, and let $P_{\rm odd}^{(d)}$ be antipodal odd projection.
Fix computable real coefficients $\tau,c_0,\ldots,c_K$ and put
\[
 T^{(d)}=\tau P_{\rm odd}^{(d)}+
                 \sum_{k=0}^Kc_kP_{2k+1}^{(d)},\qquad
 \lambda_{2k+1}=\tau+c_k\ (k\leq K).
\]
Its remaining odd eigenvalues are $\tau$.  Suppose
\begin{equation}\label{eq:new-odd-spectral-cube}
 \tau+c_0=1,\qquad |\tau|\leq1,\qquad
                    |\tau+c_k|\leq1\quad(1\leq k\leq K).
\end{equation}
Define
\[
 N(T)=\sup_{d\geq1}\sup_{|f|,|g|\leq1}
             \langle f,T^{(d)}g\rangle.
\]
The absolute value may equivalently be inserted in this definition.
Odd signs suffice: odd projection preserves the unit ball of $L_\infty$,
and a bilinear functional over bounded odd functions has the same
supremum as over odd signs.

\begin{proposition}[Exact radial Hermite energies]
\label{prop:new-all-radial-energies}
For $X\sim\gamma_d$, write $R=\|X\|$, $V=X/R$, and
\[
 a_d=\frac{(\mathbb ER)^2}{d}
 =\frac{2\Gamma((d+1)/2)^2}{d\Gamma(d/2)^2},\qquad \alpha=d/2.
\]
With the generalized Laguerre polynomial $L_k^\alpha$, one has
\begin{align}
 P_{2k+1}^{(d)}V_i
  &=\frac{\mathbb ER}{d}\frac{(1/2)_k}{(\alpha+1)_k}
             X_iL_k^\alpha(R^2/2),\label{eq:new-vector-projection}\\
 e_{2k+1}(d):=\sum_i\|P_{2k+1}^{(d)}V_i\|_2^2
  &=a_d\frac{(1/2)_k^2}{k!(\alpha+1)_k}.
             \label{eq:new-radial-energy}
\end{align}
In particular
\[
 e_1=a_d,\quad e_3=\frac{a_d}{2(d+2)},\quad
 e_5=\frac{9a_d}{8(d+2)(d+4)},\quad
 e_7=\frac{75a_d}{16(d+2)(d+4)(d+6)}.
\]
The exact natural-vector objective is
\begin{equation}\label{eq:new-exact-natural-value}
 q_d(T)=\tau+\sum_{k=0}^Kc_ke_{2k+1}(d)\longrightarrow1.
\end{equation}
\end{proposition}
\begin{proof}
Gaussian projection commutes with orthogonal rotations.  Thus the
projection of the equivariant vector field $V$ onto degree $2k+1$ is
a constant times $X_iL_k^\alpha(R^2/2)$.  Indeed an equivariant vector
polynomial is $X_i$ times a radial polynomial, and orthogonality to
the lower odd degrees is precisely Laguerre orthogonality with
weight $t^\alpha e^{-t}$.
Put $t=R^2/2$.  Radial integration and Laguerre orthogonality give
\[
 \mathbb E[X_i^2L_k^\alpha(t)^2]=\frac{(\alpha+1)_k}{k!},\qquad
 \mathbb E[V_iX_iL_k^\alpha(t)]
       =\frac{\mathbb ER}{d}\frac{(1/2)_k}{k!}.
\]
For the second identity, use
\[
 \int_0^\infty t^{\alpha-1/2}e^{-t}L_k^\alpha(t)\,dt
       =\Gamma(\alpha+1/2)\frac{(1/2)_k}{k!},
\]
which follows by integrating the generating function
$(1-z)^{-\alpha-1}\exp[-tz/(1-z)]$ and comparing coefficients.
Taking the ratio proves \eqref{eq:new-vector-projection}; its squared
norm proves \eqref{eq:new-radial-energy}.  Parseval gives
$\sum_{k\geq0}e_{2k+1}(d)=1$.  Finally $a_d\to1$ and every fixed
$e_{2k+1}(d)$ with $k\geq1$ tends to zero.  Since $\tau+c_0=1$,
\eqref{eq:new-exact-natural-value} follows.
\end{proof}

\begin{proposition}[An explicit reduced matrix sequence]
\label{prop:new-odd-tail-matrix-sequence}
For $n\geq2$, set $d=n^2$.  Partition each Gaussian coordinate into
intervals of mesh $n^{-2}$ from $-n$ to $n$, together with the two
exterior intervals.  Retain one cell $C$ from each antipodal pair,
using the cells with first coordinate negative.  Set
\[
 w_C=2\gamma_d(C),\qquad D_w=\operatorname{diag}(w_C),\qquad
 F_j[C,\xi]=2\int_C H_\xi\,d\gamma_d\quad(|\xi|=j),
\]
where $H_\xi$ is the normalized product Hermite polynomial.  Define
\begin{equation}\label{eq:new-odd-tail-seed}
 A_n(T)=\tau D_w+\sum_{k=0}^Kc_kF_{2k+1}F_{2k+1}^T.
\end{equation}
Then
\begin{equation}\label{eq:new-tail-matrix-limits}
 C(A_n(T))\to N(T),\qquad V(A_n(T))\to1,
 \qquad G(A_n(T))\to\frac1{N(T)}.
\end{equation}
Choose dyadic rational approximants with total entrywise error at most
$\eta_n=n^{-4}$.  Clear their denominators, delete zero rows and
columns, and apply the exact absolute row/column cloning reduction,
the symmetric off-diagonal dilation, and the signed balanced copy.
The resulting rational matrices $H_n(T)$ are symmetric, zero-diagonal,
have zero row and column sums, and have $|H_n(T)|$ doubly stochastic.
Their ratios have the same limit $1/N(T)$.
\end{proposition}
\begin{proof}
Odd extension of the cell values turns the finite bilinear form
\eqref{eq:new-odd-tail-seed} into $\langle f,T^{(d)}g\rangle$.
The diagonal term is exact because for odd cell functions
\[
 \langle f,P_{\rm odd}^{(d)}g\rangle
           =\sum_Cw_Cf_Cg_C.
\]
The remaining terms are the exact Hermite projection pairings.
Thus $C(A_n(T))\leq N(T)$.  Every pair of bounded odd functions in a
fixed dimension can be approximated in $L_2$ by conditional averages
on the first coordinates of these grids as $n\to\infty$.
Continuous compactly supported approximation proves this even though
the grids are not nested.  The other coordinates integrate out, and
$T^{(d)}$ is consistent under such cylindrical embeddings.  Since
its $L_2$ norm is one, the pairings converge.  Finite cube maxima
occur at sign vertices, proving $\liminf C(A_n(T))\geq N(T)$.

For vector-valued odd cell functions of pointwise norm one, the same
$L_2$ operator bound and Cauchy--Schwarz give $V(A_n(T))\leq1$.
Use the unit vector at the clipped cell midpoint as the trial field.
The quantitative estimate below proves that its objective tends to
$q_d(T)$ and hence to one.  This proves the vector limit.  Positivity
is explicit: for $s=\operatorname{sgn}X_1$, first-chaos energy is
$2/\pi$ and all remaining odd eigenvalues are at least $-1$, whence
$N(T)\geq\langle s,Ts\rangle\geq4/\pi-1>0$.

Both finite norms change by at most total entrywise approximation
error.  Consequently rational approximation preserves both limits.
For completeness, if $Z$ is the cleared integer seed, put
$r_i=\sum_j|Z_{ij}|$, $s_j=\sum_i|Z_{ij}|$, and replace its entry by
$Z_{ij}/(r_is_j)$ on each of the $r_i$ row and $s_j$ column copies.
Both norms are unchanged, because every original row and column has
total absolute copying weight one; the absolute matrix is doubly
stochastic.  Off-diagonal symmetric dilation multiplies both norms
by two.  Tensoring with
$\tfrac12\left(\begin{smallmatrix}1&-1\\-1&1\end{smallmatrix}\right)$
multiplies both by two again and imposes zero sums.  All ratios are
therefore preserved exactly.
\end{proof}

\begin{proposition}[The general asymptotic vector value]
\label{prop:new-general-vector-value}
Drop \eqref{eq:new-odd-spectral-cube} and assume only that $T$ is
nonzero.  Write
\[
 M=\max\bigl(|\tau|,|\tau+c_0|,\ldots,|\tau+c_K|\bigr).
\]
The seed \eqref{eq:new-odd-tail-seed} and its rational reductions satisfy
\[
 C(A_n(T))\to N(T),\qquad V(A_n(T))\to M,
                    \qquad G(H_n(T))\to M/N(T).
\]
For every fixed odd Hermite degree $j$ and $d\geq j$, the square-free
tensor trial field used in the proof has norm-square variance
\[
 \epsilon_{d,j}^2=
 \sum_{r=1}^j
 \frac{\binom jr\binom{d-j}{j-r}}{\binom dj}(3^r-1)
                  \leq\frac{j(3^j-1)}d.
\]
Terms with an impossible intersection size have binomial coefficient
zero.  Its unit normalization attains $|\lambda_j|$
within $2M\epsilon_{d,j}$ in continuum vector objective.
\end{proposition}
\begin{proof}
Let $m=\binom dj$ and define the vector with components indexed by
$j$-subsets $S$ of $\{1,\ldots,d\}$,
\[
 Y_S(X)=m^{-1/2}\prod_{i\in S}X_i,\qquad U=Y/\|Y\|.
\]
The exceptional zero set has Gaussian measure zero.  Every component
has pure degree $j$, and $\mathbb E\|Y\|^2=1$.  For two subsets
$S,S'$, their norm-square moment is $3^{|S\cap S'|}$.
Counting intersections proves the variance formula.  If $R$ is their
intersection size, $\mathbb ER=j^2/d$ and
$3^r-1\leq r(3^j-1)/j$ for $0\leq r\leq j$, giving its upper bound.
The elementary inequality $(\sqrt{x}-1)^2\leq(x-1)^2$ yields
\[
 \|U-Y\|_{L_2}^2=\mathbb E(\|Y\|-1)^2
                   \leq\epsilon_{d,j}^2.
\]
Thus
$|\langle U,TU\rangle-\lambda_j|\leq2M\epsilon_{d,j}$.
If $\lambda_j<0$, use $-U$ on the right.  Any selected nonzero odd
eigenvalue is covered; for the tail eigenvalue choose an odd
$j>2K+1$.  This proves that the continuum supremum of the vector
objective over all dimensions is $M$.

For each fixed trial dimension, conditional cell averages of $U$
on the first coordinates converge in $L_2$ as $n\to\infty$.
They have pointwise norm at most one; completing left and right
ball vectors in separate orthogonal coordinates gives admissible
finite unit vectors with the same cross pairings.  Therefore the
seed vector norms have $\liminf\geq M$.  The $L_2$ norm of $T$ gives
the reverse bound.  The scalar convergence and rational reductions
are the same as in the preceding proof, with operator norm $M$
replacing one.  A nonzero multiplier has $N(T)>0$, since otherwise
all its pairings against bounded functions vanish, contradicting its
nonzero Hermite action.
\end{proof}

\paragraph{An explicit finite error through degree seven.}
Here $K\leq3$.  Set missing $c_k$ equal to zero and put
\[
 \lambda_*=\min(\tau,\tau+c_0,\tau+c_1,\tau+c_2,\tau+c_3),
 \qquad C_* = |\tau-\lambda_*|.
\]
Define the following nonnegative rational expressions in $d$ and
the absolute coefficient values:
\begin{align*}
 B_0&=|c_0|+|c_1|/2+3|c_2|/8+5|c_3|/16,\\
 B_1&=\frac{|c_1|/2+3|c_2|/4+15|c_3|/16}{d+2},\\
 B_2&=\frac{3|c_2|/8+15|c_3|/16}{(d+2)(d+4)},\\
 B_3&=\frac{5|c_3|/16}{(d+2)(d+4)(d+6)}.
\end{align*}
Let $J_0=1$, $J_j=\prod_{a=1}^j(d+2a-1)$ and
$M_j=\prod_{a=0}^{j-1}(d+2a)$, with $M_0=1$.  Put
\begin{align*}
 Q_d^*&=C_*^2+2C_*\sum_{j=0}^3B_jJ_j
          +\frac1d\sum_{i,j=0}^3B_iB_jM_{i+j+1},\\
 \mathcal L_n(T)&=
 \frac{d}{n^4}\left[
 \frac{C_*}{d-2}+\frac{B_0}{d-1}+B_1
 +(d+1)B_2+(d+1)(d+3)B_3\right]
       +4\sqrt{Q_d^*}\sqrt{2n\varphi(n)},\quad d=n^2.
\end{align*}
Then $\mathcal L_n(T)=O(n^{-4})$ and
\begin{equation}\label{eq:new-tail-finite-lower}
 V(A_n(T))\geq q_{n^2}(T)-\mathcal L_n(T).
\end{equation}
In particular, if a separately proved all-sign bound is $N(T)\leq B$,
the rational reduced sequence satisfies the completely explicit bound
\[
 G(H_n(T))\geq
 \frac{\max(0,q_{n^2}(T)-\mathcal L_n(T)-n^{-4})}{B+n^{-4}}.
\]
This is a lower bound; it does not identify $N(T)$ with $B$.

\begin{proof}[Proof of the error estimate]
The radial projection formula gives $TV=t_d(R)V$, where
\[
 t_d(R)=\tau+\frac{\mathbb ER}{d}R\sum_{k=0}^3
 c_k\frac{(1/2)_k}{(d/2+1)_k}L_k^{d/2}(R^2/2).
\]
Let $v$ be the clipped-midpoint unit field and $\delta=v-V$.
Both fields are odd and unit length, so
$2\delta\cdot V=-\|\delta\|^2$.  Since the minimum odd eigenvalue
is $\lambda_*$,
\[
 \langle v,Tv\rangle\geq q_d(T)
             -\mathbb E[(t_d(R)-\lambda_*)\|\delta\|^2].
\]
The coefficientwise absolute Laguerre expansion shows
\[
 (t_d(R)-\lambda_*)_+\leq
 B_d(R):=C_*+\frac{\mathbb ER}{d}
                         \sum_{j=0}^3 B_jR^{2j+1}.
\]
On the no-clipping event normalization gives
$\|\delta\|^2\leq dn^{-4}/R^2$; on its complement it is at most
four.  The latter event has probability at most $2n\varphi(n)$.
Gaussian radial moments give
\[
 \mathbb ER^{2j}=M_j,\qquad
 \mathbb ER^{2j+1}=(\mathbb ER)J_j,\qquad
 \mathbb ER^{-1}=\frac{\mathbb ER}{d-1},\quad
 \mathbb ER^{-2}=\frac1{d-2}.
\]
Use $(\mathbb ER)^2/d=a_d\leq1$ in
$\mathbb E[B_d(R)/R^2]$ and $\mathbb EB_d(R)^2$, followed by
Cauchy--Schwarz on the clipping event.  These are exactly the two
terms of $\mathcal L_n(T)$.  Finally $B_j=O(d^{-j})$, $J_j=O(d^j)$,
$M_j=O(d^j)$, which proves the asserted rate.
\end{proof}

\paragraph{Two useful subfamilies.}
For $T=P_1-b_3P_3-b_5P_5-b_7P_7$ the exact numerator is
$a_d-b_3e_3-b_5e_5-b_7e_7$.  For the normalized identity-tail family
\[
 T_{\tau,b}=(1-\tau)P_1-bP_3+\tau P_{\rm odd},
 \qquad q_d(T_{\tau,b})=\tau+(1-\tau)a_d-be_3.
\]
In the latter family the spectral-cube conditions are
$|\tau|\leq1$ and $|\tau-b|\leq1$.  The identity tail is represented
by the diagonal in \eqref{eq:new-odd-tail-seed}, and is not discarded
at finite $n$.  Its numerator has leading loss
$(1-\tau+b)/(2d)$.  For every such family, a useful numerical endpoint
still requires a proved bound for the complete scalar norm $N(T)$.
